\section{Euler characteristics}
\label{section-euler}

\noindent
In this section we prove some elementary properties of Euler characteristics
of coherent sheaves on schemes proper over fields.

\begin{definition}
\label{definition-euler-characteristic}
Let $k$ be a field. Let $X$ be a proper scheme over $k$. Let $\mathcal{F}$
be a coherent $\mathcal{O}_X$-module. In this situation the
{\it Euler characteristic of $\mathcal{F}$} is the integer
$$
\chi(X, \mathcal{F}) = \sum\nolimits_i (-1)^i \dim_k H^i(X, \mathcal{F}).
$$
For justification of the formula see below.
\end{definition}

\noindent
In the situation of the definition only a finite number of the vector spaces
$H^i(X, \mathcal{F})$ are nonzero (Cohomology of Schemes, Lemma
\ref{coherent-lemma-quasi-coherence-higher-direct-images})
and each of these spaces is finite dimensional
(Cohomology of Schemes, Lemma
\ref{coherent-lemma-proper-over-affine-cohomology-finite}). Thus
$\chi(X, \mathcal{F}) \in \mathbf{Z}$ is well defined. Observe that
this definition depends on the field $k$ and not just on the pair
$(X, \mathcal{F})$.

\begin{lemma}
\label{lemma-euler-characteristic-additive}
Let $k$ be a field. Let $X$ be a proper scheme over $k$.
Let $0 \to \mathcal{F}_1 \to \mathcal{F}_2 \to \mathcal{F}_3 \to 0$
be a short exact sequence of coherent modules on $X$. Then
$$
\chi(X, \mathcal{F}_2) = \chi(X, \mathcal{F}_1) + \chi(X, \mathcal{F}_3)
$$
\end{lemma}

\begin{proof}
Consider the long exact sequence of cohomology
$$
0 \to H^0(X, \mathcal{F}_1) \to H^0(X, \mathcal{F}_2) \to
H^0(X, \mathcal{F}_3) \to H^1(X, \mathcal{F}_1) \to \ldots
$$
associated to the short exact sequence of the lemma. The rank-nullity theorem
in linear algebra shows that
$$
0 = \dim H^0(X, \mathcal{F}_1) - \dim H^0(X, \mathcal{F}_2)
+ \dim H^0(X, \mathcal{F}_3) - \dim H^1(X, \mathcal{F}_1) + \ldots
$$
This immediately implies the lemma.
\end{proof}

\begin{lemma}
\label{lemma-chi-tensor-finite}
Let $k$ be a field. Let $X$ be a proper scheme over $k$. Let $\mathcal{F}$
be a coherent sheaf with $\dim(\text{Supp}(\mathcal{F})) \leq 0$.
Then
\begin{enumerate}
\item $\mathcal{F}$ is generated by global sections,
\item $H^0(X, \mathcal{F}) =
\bigoplus_{x \in \text{Supp}(\mathcal{F})} \mathcal{F}_x$,
\item $H^i(X, \mathcal{F}) = 0$ for $i > 0$,
\item $\chi(X, \mathcal{F}) = \dim_k H^0(X, \mathcal{F})$, and
\item
$\chi(X, \mathcal{F} \otimes \mathcal{E}) = n\chi(X, \mathcal{F})$
for every locally free module $\mathcal{E}$ of rank $n$.
\end{enumerate}
\end{lemma}

\begin{proof}
By Cohomology of Schemes, Lemma \ref{coherent-lemma-coherent-support-closed}
we see that $\mathcal{F} = i_*\mathcal{G}$ where $i : Z \to X$ is the inclusion
of the scheme theoretic support of $\mathcal{F}$ and where $\mathcal{G}$
is a coherent $\mathcal{O}_Z$-module. By definition of the scheme theoretic
support the underlying topological space of $Z$ is $\text{Supp}(\mathcal{F})$.
Since the dimension of $Z$ is $0$, we see $Z$ is affine (Properties, Lemma
\ref{properties-lemma-locally-Noetherian-dimension-0}).
Hence $\mathcal{G}$ is globally generated and the higher
cohomology groups of $\mathcal{G}$ are zero
(Cohomology of Schemes, Lemma
\ref{coherent-lemma-quasi-coherent-affine-cohomology-zero}).
In fact, by Lemma \ref{lemma-algebraic-scheme-dim-0}
the scheme $Z$ is a finite disjoint union of spectra
of local Artinian rings. Thus correspondingly
$H^0(Z, \mathcal{G}) = \bigoplus_{z \in Z} \mathcal{G}_z$.
The cohomologies of $\mathcal{F}$ and $\mathcal{G}$ agree by
Cohomology of Schemes, Lemma \ref{coherent-lemma-relative-affine-cohomology}.
Thus $H^i(X, \mathcal{F}) = 0$ for $i > 0$ and
$H^0(X, \mathcal{F}) = H^0(Z, \mathcal{G})$.
In particular we have (3) is true.
For $z \in Z$ corresponding to $x \in \text{Supp}(\mathcal{F})$
we have $\mathcal{G}_z = (i_*\mathcal{G})_x = \mathcal{F}_x$.
We conclude that (2) holds. Of course (2) implies (1).
We have (4) by definition of the Euler characteristic $\chi(X, \mathcal{F})$
and (3). By the projection formula
(Cohomology, Lemma \ref{cohomology-lemma-projection-formula}) we have
$$
i_*(\mathcal{G} \otimes i^*\mathcal{E}) = \mathcal{F} \otimes \mathcal{E}.
$$
Since $Z$ has dimension $0$ the locally free sheaf $i^*\mathcal{E}$
is isomorphic to $\mathcal{O}_Z^{\oplus n}$ and arguing as above
we see that (5) holds.
\end{proof}

\begin{lemma}
\label{lemma-euler-characteristic-extend-base-field}
Let $k'/k$ be an extension of fields. Let $X$ be a proper scheme
over $k$. Let $\mathcal{F}$ be a coherent sheaf on $X$.
Let $\mathcal{F}'$ be the pullback of $\mathcal{F}$ to $X_{k'}$.
Then $\chi(X, \mathcal{F}) = \chi(X', \mathcal{F}')$.
\end{lemma}

\begin{proof}
This is true because
$$
H^i(X_{k'}, \mathcal{F}') = H^i(X, \mathcal{F}) \otimes_k k'
$$
by flat base change, see
Cohomology of Schemes, Lemma \ref{coherent-lemma-flat-base-change-cohomology}.
\end{proof}

\begin{lemma}
\label{lemma-euler-characteristic-morphism}
Let $k$ be a field. Let $f : Y \to X$ be a morphism of proper schemes over
$k$. Let $\mathcal{G}$ be a coherent $\mathcal{O}_Y$-module. Then
$$
\chi(Y, \mathcal{G}) = \sum (-1)^i \chi(X, R^if_*\mathcal{G})
$$
\end{lemma}

\begin{proof}
The formula makes sense: the sheaves $R^if_*\mathcal{G}$ are coherent
and only a finite number of them are nonzero, see
Cohomology of Schemes,
Proposition \ref{coherent-proposition-proper-pushforward-coherent} and
Lemma \ref{coherent-lemma-quasi-coherence-higher-direct-images}.
By Cohomology, Lemma \ref{cohomology-lemma-Leray} there is a spectral
sequence with
$$
E_2^{p, q} = H^p(X, R^qf_*\mathcal{G})
$$
converging to $H^{p + q}(Y, \mathcal{G})$. By finiteness of cohomology
on $X$ we see that only a finite number of $E_2^{p, q}$ are nonzero
and each $E_2^{p, q}$ is a finite dimensional vector space. It follows
that the same is true for $E_r^{p, q}$ for $r \geq 2$ and that
$$
\sum (-1)^{p + q} \dim_k E_r^{p, q}
$$
is independent of $r$. Since for $r$ large enough we have
$E_r^{p, q} = E_\infty^{p, q}$ and since convergence means there
is a filtration on $H^n(Y, \mathcal{G})$ whose graded pieces are
$E_\infty^{p, q}$ with $p + q = n$ (this is the meaning of convergence
of the spectral sequence), we conclude.
Compare also with the more general
Homology, Lemma \ref{homology-lemma-biregular-ss-relation-in-K0}.
\end{proof}








